//323111047

\section{Introducción}

Planteamiento del problema. El reto técnico de este proyecto consiste en desarrollar una aplicación de línea de comandos en lenguaje C que gestione un carrito de compra mediante el uso de estructuras de datos dinámicas. El problema central radica en implementar la bidireccionalidad en la navegación de los productos, permitiendo que el usuario visualice la lista tanto en el orden de ingreso como en orden inverso. Para lograr esto, es indispensable diseñar un nodo que contenga los atributos de nombre, identificador y precio, además de gestionar punteros dobles que permitan el recorrido en ambos sentidos y la eliminación de elementos específicos sin comprometer la integridad de la memoria dinámica.

\vspace{0.4cm}

Motivación. La importancia de resolver este problema se basa en la necesidad de diseñar sistemas de gestión de inventarios y ventas que sean altamente flexibles y eficientes. En el comercio electrónico real, un carrito de compras debe permitir al usuario revisar sus selecciones de forma ágil y realizar modificaciones constantes, como eliminar artículos específicos, sin que el sistema se vuelva inestable. Estudiar estas estructuras dinámicas permite al estudiante comprender cómo optimizar el acceso a la información y cómo administrar recursos computacionales de manera profesional, preparándolo para el desarrollo de aplicaciones comerciales de alto rendimiento.

\vspace{0.4cm}

Objetivos. Se busca que el alumno demuestre dominio en las funciones de inserción, eliminación y recorrido bidireccional de nodos, validando que el flujo de datos sea consistente en ambos sentidos de la lista. Finalmente, se pretende verificar que el programa genere un ticket de salida preciso con el desglose de productos y la sumatoria total del costo total.